% Theorem 3.1.2* (CORRECTED): all ideals of norm p^m in O_d, p | f.
% Corrected p-uniform form per ERRATA E1, E2, E4. Do NOT restore the thesis's
% separate five-case 2-adic display: two of its branches are false as stated.

\begin{theorem}[Ideals of norm $p^m$ in $\mathcal{O}_d$, $p \mid f$; corrected Theorem 3.1.2]
\label{thm-3-1-2}
\lean{QuadraticOrder.idealCount_prime_pow_of_le, QuadraticOrder.idealCount_prime_pow_ramified, QuadraticOrder.idealCount_prime_pow_inert, QuadraticOrder.idealCount_prime_pow_split}
\uses{notation, infra-index}
Let $D < 0$ be a fundamental discriminant, $K = \mathbb{Q}(\sqrt D)$ with maximal
order $\mathcal{O}_K = \mathbb{Z}[(D + \sqrt D)/2]$, let $f \geq 1$ and
$d := D f^2$, and let $\mathcal{O}_d = \mathbb{Z}[\tau]$, $\tau := (d + \sqrt d)/2$,
be the order of discriminant $d$ and conductor $f$, so that
$\mathcal{O}_d \cong \mathbb{Z}[x]/(g)$ with
$g(x) := x^2 - d\,x + \tfrac{d^2 - d}{4}$. Write
$[\mathcal{O}_d : \mathfrak a]$ for the (finite) index of a nonzero ideal
$\mathfrak a \subseteq \mathcal{O}_d$, i.e.\ its norm.

Let $p$ be a prime with $p \mid f$, set $w := v_p(f) \geq 1$, let $m \geq 0$, and
let $\varepsilon := 1$ if $m$ is even and $\varepsilon := 0$ if $m$ is odd. Define
\[
v := v_p(d/4) :=
\begin{cases}
v_p(d) & p \text{ odd},\\
v_2(d) - 2 & p = 2 \quad (\text{legitimate: } 2 \mid f \Rightarrow 4 \mid d),
\end{cases}
\qquad\text{so } v =
\begin{cases}
v_p(D) + 2w & p \text{ odd},\\
v_2(D) + 2w - 2 & p = 2.
\end{cases}
\]
Then the number of ideals of norm $p^m$ in $\mathcal{O}_d$ (invertible or not) is
given as follows.
\begin{itemize}
\item If $v \geq m$:
\[
\#\{\mathfrak a \subseteq \mathcal{O}_d : [\mathcal{O}_d : \mathfrak a] = p^m\}
= \sum_{i=0}^{\lfloor m/2 \rfloor} p^i
= \frac{p^{\lfloor m/2 \rfloor + 1} - 1}{p - 1}.
\]
\item If $v < m$:
\[
\#\{\mathfrak a \subseteq \mathcal{O}_d : [\mathcal{O}_d : \mathfrak a] = p^m\}
= \begin{cases}
\displaystyle\sum_{i=0}^{w} p^i
  & p \text{ ramified in } \mathcal{O}_K \ (p \mid D),\\[2ex]
\displaystyle\sum_{i=0}^{w + \varepsilon - 1} p^i
  & p \text{ inert in } \mathcal{O}_K,\\[2ex]
\displaystyle\sum_{i=0}^{w + \varepsilon - 1} p^i + p^w\,(m + 1 - 2w - \varepsilon)
  & p \text{ split in } \mathcal{O}_K.
\end{cases}
\]
\end{itemize}
The statement is uniform in $p$: there is no separate display for $p = 2$.
\end{theorem}

\begin{proof}
\uses{lem-3-2-4, lem-3-2-6}
Write $r(p^m) := \#\{\mathfrak a \subseteq \mathcal{O}_d :
[\mathcal{O}_d : \mathfrak a] = p^m\}$ and $\mu := m \bmod 2 = 1 - \varepsilon$.

\emph{Step 1 (reduction to root counting).} By the normal-form bijection of
Lemma~\ref{lem-3-2-4} (in its corrected form, which applies to \emph{all} ideals of
index $p^m$ --- since $p \mid f$, the polynomial $g \bmod p$ is the square of a linear
polynomial, so $\mathcal{O}_d/p\mathcal{O}_d$ is local, there is a unique prime of
$\mathcal{O}_d$ above $p$, and every ideal of $p$-power index is primary for it ---
and which includes the minimal-integer normalization
$\mathfrak a_{k,A} \cap \mathbb{Z} = p^k\mathbb{Z}$ preventing double counting), the
ideals of index $p^m$ are exactly the
\[
\mathfrak a_{k,A} = p^k\,\mathbb{Z} \oplus p^{m-k}(\tau - A)\,\mathbb{Z},
\qquad
\lceil m/2 \rceil \leq k \leq m,\quad
[A] \in \mathbb{Z}/p^{2k-m},\quad
g(A) \equiv 0 \bmod p^{2k-m},
\]
each exactly once. Hence
\[
r(p^m) = \sum_{k=\lceil m/2 \rceil}^{m} N_k,
\qquad
N_k := \#\{A \in \mathbb{Z}/p^{2k-m} : g(A) \equiv 0 \bmod p^{2k-m}\}.
\]
Set $s := 2k - m$; as $k$ runs over its range, $s$ runs over the integers of the same
parity as $m$ with $\mu \leq s \leq m$, and
$k - \lceil m/2 \rceil = (s - \mu)/2 = \lfloor s/2 \rfloor$.

\emph{Step 2 (completing the square).} In $\mathbb{Z}[\tfrac12][x]$ one has
$g(x) = (x - d/2)^2 - d/4$. If $p$ is odd, $2$ is invertible mod $p^s$ and the
substitution $x \mapsto x + d\,\bar 2$ (with $\bar 2 := 2^{-1} \bmod p^s$) is a
bijection of $\mathbb{Z}/p^s$; if $p = 2$ then $4 \mid d$ and the identity is exact
in $\mathbb{Z}[x]$. In both cases
\[
N_k = \#\{y \in \mathbb{Z}/p^s : y^2 \equiv q \bmod p^s\},
\qquad
q := \begin{cases} d\,\bar 4 & p \text{ odd},\\ d/4 & p = 2,\end{cases}
\]
and any integer lift of $q$ has $p$-adic valuation exactly $v$ when $v < s$, while
$q \equiv 0 \bmod p^s$ when $v \geq s$.

\emph{Step 3 (valuation dichotomy; the case $v \geq m$).} Put
$\kappa := (v + m)/2$, so $k \leq \kappa \iff s \leq v$. If $s \leq v$ then
$y^2 \equiv 0 \bmod p^s$, whose solutions are exactly the multiples of
$p^{\lceil s/2 \rceil}$, giving
$N_k = p^{\lfloor s/2 \rfloor} = p^{\,k - \lceil m/2 \rceil}$. If $v \geq m$ every
$k$ in range satisfies $s \leq m \leq v$, so
\[
r(p^m) = \sum_{k = \lceil m/2 \rceil}^{m} p^{\,k - \lceil m/2 \rceil}
= \sum_{i=0}^{\lfloor m/2 \rfloor} p^i,
\]
proving the first branch (this also covers $m = 0$). Assume $v < m$ from now on.

\emph{Step 4 (the sub-$\kappa$ geometric block).} The terms with $k \leq \kappa$
total $\Sigma_0 := \sum_{i=0}^{i_0} p^i$ with
$i_0 := \lfloor (v+m)/2 \rfloor - \lceil m/2 \rceil$ (an empty sum, $= 0$, if
$i_0 < 0$). A direct computation gives: $i_0 = (v-1)/2$ if $v$ is odd (both parities
of $m$), and $i_0 = v/2 - 1 + \varepsilon$ if $v$ is even.

\emph{Step 5 (beyond $\kappa$).} The remaining terms have
$s \in S := \{s : v < s \leq m,\ s \equiv m \bmod 2\}$, which is nonempty; its least
element $s_1$ is $v+1$ or $v+2$.

\emph{(a) $v$ odd} --- exactly the cases $p$ odd with $p \mid D$ (then $v_p(D) = 1$,
$v = 2w+1$) and $p = 2$ with $D \equiv 8 \bmod 16$ ($v = 2w+1$); both ramified. Then
$N_k = 0$ for all $s > v$: a solution would force $v_p(y^2) = v$ by the ultrametric
inequality (any integer lift of $q$ has valuation $v < s$), impossible since
$v_p(y^2)$ is even. Hence $r(p^m) = \Sigma_0 = \sum_{i=0}^{w} p^i$, the ramified
formula.

\emph{(b) $v$ even.} Write $r := v/2$ and choose an integer $u$ coprime to $p$ with
$q \equiv p^{2r} u \bmod p^s$. Concretely: for $p$ odd unramified, $r = w$ and
$u \equiv D (f/p^w)^2\, \bar 4 \bmod p$, so the Legendre symbol satisfies
$(\tfrac up) = (\tfrac Dp)$; for $p = 2$ with $D$ odd, $r = w - 1$ and, writing
$f = 2^w f'$ with $f'$ odd, $u = D f'^2 \equiv D \bmod 8$; for $p = 2$ with
$D \equiv 4 \bmod 8$, $r = w$ and $u = (D/4) f'^2 \equiv 3 \bmod 4$. Since
$2r = v < s$ on $S$, Lemma~\ref{lem-3-2-6} (corrected form) applies: for $p$ odd
there are $2p^r$ solutions if $(\tfrac up) = 1$ and none otherwise; for $p = 2$
there are $2^r$ solutions if $s - 2r = 1$ (unconditionally), $2^{r+1}$ if
$s - 2r = 2$ and $u \equiv 1 \bmod 4$, $2^{r+2}$ if $s - 2r \geq 3$ and
$u \equiv 1 \bmod 8$, and none otherwise. \textbf{The case distinction is governed by
$s - 2r$, not by $s$}; failing to track this is precisely the error in the thesis
proof (ERRATA E3).

\emph{(b1) $p$ odd inert} ($(\tfrac Dp) = -1$): no contributions, and
$r(p^m) = \Sigma_0 = \sum_{i=0}^{w-1+\varepsilon} p^i$, the inert formula
($v = 2w$).

\emph{(b2) $p$ odd split} ($(\tfrac Dp) = 1$): every $s \in S$ contributes $2p^w$.
Counting, $\#S = (m + 1 - 2w - \varepsilon)/2 \geq 1$, so
\[
r(p^m) = \sum_{i=0}^{w-1+\varepsilon} p^i
+ 2p^w \cdot \frac{m + 1 - 2w - \varepsilon}{2}
= \sum_{i=0}^{w+\varepsilon-1} p^i + p^w (m + 1 - 2w - \varepsilon),
\]
the split formula.

\emph{(b3) $p = 2$, $D \equiv 4 \bmod 8$} (ramified, $v = 2w$, $r = w$,
$u \equiv 3 \bmod 4$): the only admissible value is $s - 2r = 1$, i.e.\
$s = 2w + 1 \in S$ iff $m$ is odd (and then $2w + 1 \leq m$ since $v = 2w < m$),
contributing $2^w$; the cases $s - 2r = 2$ and $s - 2r \geq 3$ are excluded by
$u \equiv 3 \bmod 4$. Hence
\[
r(2^m) = \sum_{i=0}^{w-1+\varepsilon} 2^i + (1 - \varepsilon)2^w
= \sum_{i=0}^{w} 2^i
\]
for both parities of $m$ (for $m$ even this is $\Sigma_0$ alone; for $m$ odd it is
$(2^w - 1) + 2^w$), the ramified formula. This corrects ERRATA E2: the thesis
claimed $2^{w+\varepsilon} - 1$, wrong for $m$ odd.

\emph{(b4) $p = 2$, $D$ odd} ($v = 2w - 2$, $r = w - 1$, $u \equiv D \bmod 8$; here
$s - 2r = s - 2w + 2$). By Lemma~\ref{lem-3-2-6}: $s = 2w - 1$ ($\in S$ iff $m$ odd)
contributes $2^{w-1}$ unconditionally; $s = 2w$ ($\in S$ iff $m$ even) contributes
$2^w$ since $u \equiv D \equiv 1 \bmod 4$; and each $s \geq 2w + 1$ in $S$
contributes $2^{w+1}$ if $D \equiv 1 \bmod 8$ ($2$ split) and $0$ if
$D \equiv 5 \bmod 8$ ($2$ inert). In particular the least element $s_1$ of $S$
contributes exactly $2^{\,w - 1 + \varepsilon}$ \emph{for every} $w \geq 1$ --- the
term the thesis found only at $w = 1$ (ERRATA E1/E3). Therefore, with
$\Sigma_0 = \sum_{i=0}^{w-2+\varepsilon} 2^i = 2^{w-1+\varepsilon} - 1$ (empty when
$w = 1$, $m$ odd):
\[
\text{inert:}\quad
r(2^m) = (2^{w-1+\varepsilon} - 1) + 2^{w-1+\varepsilon}
= \sum_{i=0}^{w+\varepsilon-1} 2^i;
\]
\[
\text{split:}\quad
r(2^m) = (2^{w-1+\varepsilon} - 1) + 2^{w-1+\varepsilon}
+ 2^{w+1}\cdot\frac{m+1-2w-\varepsilon}{2}
= \sum_{i=0}^{w+\varepsilon-1} 2^i + 2^w(m + 1 - 2w - \varepsilon),
\]
where the number of $s \in S$ with $s \geq 2w+1$ is
$(m + 1 - 2w - \varepsilon)/2 \geq 0$, exactly as in (b2).

Cases (a) and (b3) together are exactly ``$p$ ramified''; (b1) and the inert half of
(b4) are ``$p$ inert''; (b2) and the split half of (b4) are ``$p$ split''. Every
total matches the claimed $p$-uniform display.
\end{proof}

\paragraph{Divergence from the thesis.} This node states the corrected form of
Theorem 3.1.2 (ERRATA E1, E2, E4): the thesis's $2$-adic inert branch for
$w \neq 1$ ($2^{w-1+\varepsilon} - 1$) is replaced by $2^{w+\varepsilon} - 1$ for all
$w \geq 1$ (E1); the $2$-adic ramified branch with $D \equiv 4 \bmod 8$
($2^{w+\varepsilon} - 1$) is replaced by $2^{w+1} - 1$ for both parities of $m$ (E2);
after these corrections the five-case $2$-adic display coincides with the odd-$p$
formulas at $p = 2$ and is therefore dissolved into the single $p$-uniform statement
above (E4). Sanity anchors (verified by explicit enumeration; ERRATA E1/E2):
$d = -16$, $m = 3 \mapsto 3$; $d = -48$, $m = 3 \mapsto 3$ and $m = 4 \mapsto 7$;
$d = -192$, $m = 5 \mapsto 7$. These are enforced at build time as
\texttt{native\_decide} regression vectors in
\texttt{QuadraticOrder/Harness/Regression.lean}.
